% 00_overview.tex - what the workshop covers, what you need, how to read it
\begin{goal}
\textbf{What you do.} You take an ordinary Raven model, Tutorial Creek (a small catchment with two subbasins), and give it
a groundwater model that Raven builds. You write aquifer classes and profiles. You read a DEM and a bedrock surface, as GeoTIFF
and as ESRI ASCII grids, and give the projection of the input files (\cmd{:InputCRS}). You choose the grid and refine it.
You add head boundaries (with and without a \cmd{:BoundaryHead} series), a pumping well and observation wells. Then you run the
model and check its water budget.

\textbf{What you need.} A Linux computer (x86\_64) with a terminal, or Windows with WSL. Python~3 with numpy, pandas and
matplotlib (Step~\ref{st:python}). A text editor, for example \cmd{nano}. Workshop~1 is helpful but not needed.

\textbf{Time.} About one hour.
\end{goal}
